%% kapa: exam=1 round=27 year=2016 period=2 form=A subject=appraisal_law subject_ko=감정평가관계법규 seq=7 number=7
%% source: Q-Net 감정평가사 기출문제 (cst00309) artlSeq=5169402 file=extracted/1008/5169402/제27회 감정평가사 1차 문제/제27회 감정평가사 2교시문제(A형).pdf
%% conversion: kordoc 4.24.51 -> tools/gen1.py (auto); review=proofread-sonnet; flags=none
\begin{question}{7}
국토의 계획 및 이용에 관한 법령상 지구단위계획에 관한 설명으로 옳은 것은?
\begin{choices}
  \item 지구단위계획은 도시ㆍ군기본계획으로 결정한다.
  \item 지구단위계획의 수립기준 등은 시ㆍ도지사가 정한다.
  \item 지구단위계획구역은 계획관리지역에 한하여 지정할 수 있다.
  \item 계획관리지역 내에 지정하는 지구단위계획구역에 대해서는 당해 지역에 적용되는 건폐율의 200퍼센트 및 용적률의 150퍼센트 이내에서 완화하여 적용할 수 있다.
  \item 용도지역을 변경하는 지구단위계획에는 건축물의 용도제한이 반드시 포함되어야 한다.
\end{choices}
\end{question}
\answer{5}
